\subsection{Replication Ring}
\label{subsec:replication-ring}

To ensure availability, each key needs to be replicated across multiple nodes
while satisfying the following requirements:

\begin{enumerate}
    \item Each key is replicated on $N$ nodes, each from a distinct AZ.
    \item $\mathbb{E}[w_i] \leq s_i$ for each node $i$, where $w_i$ is the
        workload on node $i$ and $s_i$ is the throughput allocated by that node
        for the table, both measured in requests per second (\unit{\rps}).
\end{enumerate}

The expected workload on the $i$-th node is
\[
    \mathbb{E}[w_i] = \sum_{k \in K_i} p_k \cdot B
\]
where $K_i$ is the set of keys assigned to node $i$, $p_k$ is the
probability that key $k$ is requested and $B$ is the total throughput of the
table, in requests per second.

An extreme request distribution prevents scaling: if some key $k$ has $p_k = 1$
and is placed on node $i$, then $\mathbb[w_i] = B$, which would require
$s_i \geq B$, i.e., node $i$ must handle the entire table's workload.
To avoid this problem, we assume that keys are uniformly distributed, i.e.,
$p_k = \frac{1}{|K|}$ for each key $k$ in the table, where $K$ is the set of
keys in the table.

Under this assumption, the expected workload on node $i$ is
$\mathbb{E}[w_i] = \frac{|K_i|}{|K|} \cdot B$, so we can rewrite the second
requirement as
\[
    \frac{|K_i|}{|K|} \leq \frac{s_i}{B}.
\]

Since we don't want to waste resources, we strengthen the requirement and
enforce that $\frac{|K_i|}{|K|} = \frac{s_i}{B}$.
Summing over all nodes from the same AZ $z$, we get $\sum_{i :
\mathrm{az}(i) = z} \frac{|K_i|}{|K|} = \sum_{i : \mathrm{az}(i) = z}
\frac{s_i}{B}$, since each key is assigned to at most one node from the same
AZ, we have $\sum_{i : \mathrm{az}(i) = z} |K_i| \leq |K|$, hence
\[
    \sum_{i : \mathrm{az}(i) = z} s_i \leq B
\]
and $s_i \leq B$ for each node $i$.

Define $S = \sum s_i$, the total throughput allocated for the table across all
nodes.
By summing the second requirement over all nodes, we obtain $\frac{1}{|K|} \sum
|K_i| = \frac{S}{B}$.
Since each key is assigned to exactly $N$ nodes, we have $\sum |K_i| = N
\cdot |K|$, and so $S = B \cdot N$.

To solve this problem we will use the \emph{replication ring}.
This structure will require that the values $S$, $B$ and all $s_i$ be integers
with a granularity of $1$ request per second.
We define the \emph{replication ring} as a circular array of size $S$ where for
each position in the ring we assign a node $i$ such that:

\begin{enumerate}
    \item Every consecutive interval of length $N$ in the ring contains
        distinct nodes with different AZs.
    \item Each node $i$ occupies exactly $s_i$ positions in the ring.
\end{enumerate}

Given a key $k$, we find its starting position $p$ in the ring using its hash
value (using \texttt{XXH3}\footnote{\url{https://xxhash.com}}, a fast
non-cryptographic hash algorithm) and taking the result modulo $S$.
The key is then assigned to the $N$ nodes immediately following the
computed position, i.e., to the nodes at positions $p, p + 1, \dots, p + N - 1$
(wrapping around the ring if necessary).

Since each interval of length $N$ contains distinct nodes, each key is assigned
to $N$ nodes with different AZs, satisfying the first requirement.
Furthermore, the probability that a key is assigned to position $p$ in the ring
is $\frac{N}{S} = \frac{1}{B}$, so the expected number of keys assigned to the
position is $|K| \cdot \frac{1}{B}$.
Since each node $i$ occupies exactly $s_i$ positions in the ring and each key
is assigned at most once to the same node, the expected number of keys assigned
to node $i$ is $|K_i| = |K| \cdot \frac{s_i}{B}$, thus satisfying the second
requirement.

\subsubsection{Explicit Construction}
\label{subsubsec:explicit-construction}

For each AZ $z$, define
\[
    S_z = \sum_{i : \mathrm{az}(i) = z} s_i
\]
An AZ $z$ is \emph{saturated} if $S_z = B$.

\paragraph{Node order.}
We order nodes by first grouping them by AZ and placing all saturated AZs
before the others (the relative order of groups with the same predicate is
arbitrary).
Within each group, nodes appear in any fixed order.

\paragraph{Position order.}
We enumerate ring positions $0,\dots,S-1$ by residue classes modulo $N$: first
all multiples of $N$ in increasing order, then all positions congruent to
$1\bmod N$, and so on up to $N{-}1$.

\paragraph{Assignment.}
Traverse the nodes in the \emph{Node order} and, for each node assign it to
$s_i$ consecutive slots of the \emph{Position order}.

\paragraph{Correctness.}
Because nodes from the same AZ are consecutive in the \emph{Node order}, the
$S_z$ slots of AZ $z$ occupy a contiguous run in the \emph{Position order}.
Since each residue class contributes exactly $B$ positions and $S_z \le B$,
that run spans at most two adjacent residue classes.
Any two positions within the same residue class differ by a multiple of $N$.
If the run crosses a class boundary, the last position of class $r$ and the
first of class $r + 1$ are separated on the ring by $N + 1$ steps.
Since $S_z \le B$, the run cannot complete a full circle and wrap around to the
first position of class $r$; saturated AZs are an exception as they complete
one full circle, but all their positions belong to the same residue class
because they are placed first in the \emph{Node order}.
Hence, any two positions belonging to the same AZ are at distance at least $N$,
so every window of $N$ consecutive positions contains at most one position from
any AZ.

The construction can be expressed in pseudocode as follows in
Listing~\ref{lst:explicit-ring}.

The algorithm runs in $O(S)$ time and uses $O(S)$ space.

\begin{lstlisting}[
    language=Python,
    float=htbp,
    caption={Explicit ring construction},
    label={lst:explicit-ring}
]
def build_ring(table):
  # Positions order
  positions = []
  for x in range(table.N):
    for y in range(table.B):
      positions.append(x + y * table.N)

  # Nodes order
  nodes = table.nodes
  # Sort nodes to define the Node order:
  #   (1) saturated AZs first,
  #   (2) then by AZ identifier,
  #   (3) then by node identifier.
  nodes.sort(key=node_comparator)

  # Assignment
  ring = [None] * table.S
  idx = 0
  for node in nodes:
    for _ in range(table.s[node]):
      ring[positions[idx]] = node
      idx += 1

  return ring
\end{lstlisting}

\paragraph{Example}
Consider a table with parameters
\[
    N = 3,\quad
    B = 4,\quad
    S = B \cdot N = 12,
\]
and five nodes, each in a distinct AZ ($az(i) = i$), with allocated throughput
\[
    s_0 = 4,\; s_1 = 3,\; s_2 = 2,\; s_3 = 1,\; s_4 = 2.
\]
Since $S_{az(0)} = s_0 = 4 = B$, the AZ of node $0$ is saturated; all other AZs are non-saturated.

First we compute the \emph{Node order}.
Node $0$ takes the first position since its AZ is the only saturated one.
The remaining nodes are ordered by their AZ, so the order is:
\[
    \langle 0,\,1,\,2,\,3,\,4 \rangle.
\]

Now we compute the \emph{Position order}.
We list all positions for each residue class modulo $N$:
\[
    \begin{aligned}
        &\text{class }0: &&\langle 0,3,6,9 \rangle \\
        &\text{class }1: &&\langle 1,4,7,10 \rangle \\
        &\text{class }2: &&\langle 2,5,8,11 \rangle \\
    \end{aligned}
\]
Combining them we obtain:
\[
    P \;=\; \langle 0,3,6,9,\; 1,4,7,10,\; 2,5,8,11\rangle.
\]

Finally, we assign each node $i$ to $s_i$ consecutive positions from $P$:
\[
    \begin{aligned}
        &\text{node }0:\; P[0..3]  &&= \langle 0,3,6,9\rangle \\
        &\text{node }1:\; P[4..6]  &&= \langle 1,4,7\rangle \\
        &\text{node }2:\; P[7..8]  &&= \langle 10,2\rangle \\
        &\text{node }3:\; P[9..9]  &&= \langle 5\rangle \\
        &\text{node }4:\; P[10..11]&&= \langle 8,11\rangle \\
    \end{aligned}
\]
Thus the ring is:
\[
    \begin{array}{c|cccccccccccc}
        p   & 0 & 1 & 2 & 3 & 4 & 5 & 6 & 7 & 8 & 9 & 10 & 11 \\ \hline
        r_p & 0 & 1 & 2 & 0 & 1 & 3 & 0 & 1 & 4 & 0 & 2  & 4
    \end{array}
\]

Every window of length $N = 3$ contains nodes from distinct AZs, and each node
$i$ occupies exactly $s_i$ positions.
To place a key $k$, compute the starting position $p = \text{hash}(k) \bmod S$
and select nodes $r_p, r_{p+1}, \dots, r_{p+N-1}$ (indices modulo $S$), e.g.,
if $p = 2$ then the key is assigned to nodes $2$, $0$, and $1$
(Figure~\ref{fig:ring-example}).

\begin{figure}
    \includegraphics[width=\columnwidth]{04-implementation/ring-example}
    \caption{Replication ring example. A key hashed to position $p = 2$ is
    assigned to nodes $2$, $0$, and $1$.}
    \label{fig:ring-example}
\end{figure}

\subsubsection{Implicit Construction}
\label{subsubsec:implicit-construction}

Since the allocated throughput can be considerably high, constructing the ring
explicitly may be very expensive.
To avoid this, we look for an algorithm that directly computes the node
assigned to a position $p$ in the ring.

Let $o_p$ be the position of $p$ in the \emph{Position order}.
We find which residue class $p$ belongs to by computing $p \bmod N$; the index
within that class is given by $\left\lfloor \frac{p}{N} \right\rfloor$.
The value $p$ appears after all elements of previous classes, each containing
$B$ elements, and after $\left\lfloor \frac{p}{N} \right\rfloor$ elements of
its own class.
Hence, the index of $p$ in the \emph{Position order} is
\[
    o_p = \Bigl\lfloor \frac{p}{N} \Bigr\rfloor + (p \bmod N) \cdot B.
\]
Now we need to find the $o_p$-th assigned node, i.e., the node $i$ such that
$\sum_{j=0}^{i-1} s_j \leq o_p < \sum_{j=0}^{i} s_j$, this is the node assigned
to position $p$ in the ring.

In pseudocode, the algorithm to find the node assigned to a position $p$ in the
ring can be represented as in Listing~\ref{lst:implicit-ring}.

\begin{lstlisting}[
    language=Python,
    float=htbp,
    caption={Implicit ring construction},
    label={lst:implicit-ring}
]
def get_node_at_position(table, p):
    N = table.N
    B = table.B
    o_p = floor(p / N) + (p % N) * B

    acc = 0
    for node in range(table.nodes):
        acc += table.s[node]
        if acc > o_p:
            return node
\end{lstlisting}

We can therefore compute the node assigned to a position $p$ in the ring in
time $O(M)$ and constant memory, without explicitly building the ring.

The algorithm can be further optimized using a binary search on the prefix sums
of $s_i$, assuming they are already computed and stored.
In this case the algorithm runs in $O(\log M)$ time and uses $O(M)$ memory.

\paragraph{Example}
Let's apply the implicit construction to the example in the previous section.
We start by computing the prefix sums $c_i = \sum_{j=0}^{i-1} s_j$, obtaining
\[
    \begin{array}{c|cccccc}
        i   & 0 & 1 & 2 & 3 &  4 &  5 \\ \hline
        s_i & 4 & 3 & 2 & 1 &  2 &    \\ \hline
        c_i & 0 & 4 & 7 & 9 & 10 & 12
    \end{array}
\]

We then compute
\[
    o_p = \Bigl\lfloor \frac{p}{N} \Bigr\rfloor + (p \bmod N) \cdot B.
\]
For example, for $p = 5$ we have
\[
    o_5 = \Bigl\lfloor \tfrac{5}{3} \Bigr\rfloor + (5 \bmod 3) \cdot 4 =
    1 + 2\cdot4 = 9.
\]

We then look for $i$ such that $c_i \leq o_5 < c_{i+1}$, finding $c_3 = 9 \leq
9 < 10 = c_4$, hence the node at $p = 5$ is node~$3$.

In this way we can compute $r_p$ for any $p$ without paying the cost of
building the entire ring.

Repeating the algorithm for all $p = 0, \dots, 11$ we obtain:
\[
    \begin{array}{c|cccccccccccc}
        p   & 0 & 1 & 2 & 3 & 4 & 5 & 6 & 7 &  8 & 9 & 10 & 11 \\ \hline
        o_p & 0 & 4 & 8 & 1 & 5 & 9 & 2 & 6 & 10 & 3 &  7 & 11 \\ \hline
        r_p & 0 & 1 & 2 & 0 & 1 & 3 & 0 & 1 &  4 & 0 &  2 &  4
    \end{array}
\]

This matches the result of the explicit construction.

\subsubsection{Standard Deviation of Workloads}
\label{subsubsec:stddev}

The replication ring distributes keys to nodes proportionally to their
allocated throughput, such that the expected workload on each node is within its
capacity.
However, due to the probabilistic nature of key assignment, the actual workload
may deviate from the expected value.
We are thus interested in quantifying the standard deviation of the workload on
each node.

Assuming that the hash function distributes the keys uniformly and
independently on the ring, the probability that a random key includes node $i$
is
\[
  p_i \;=\; \frac{s_i}{B}.
\]
The number of distinct keys mapped to node $i$, denoted $|K_i|$, is
\[
  |K_i| \sim \mathrm{Binomial}\!\left(|K|,\,p_i\right).
\]
Hence
\[
  \mathbb{E}[|K_i|] = |K| \cdot p_i,
  \qquad
  \sigma_{|K_i|} = \sqrt{|K| \cdot p_i (1 - p_i)}.
\]
Since the workload on node $i$ is $w_i = |K_i| \cdot \frac{B}{|K|}$, we have
\[
    \mathbb{E}[w_i] = s_i
    \quad
    \sigma_{w_i} = \frac{B}{\sqrt{|K|}} \cdot \sqrt{p_i \cdot (1 - p_i)}.
\]

If the workload exceeds the expected value, the node will become overloaded
leading to a slowdown.
Assuming the actual workload on node $i$ to be
$\mathbb{E}[w_i] + \delta_i \sigma_{w_i}$, we can compute the slowdown factor as
\[
    f_i = \frac{\mathbb{E}[w_i] + \delta_i \sigma_{w_i}}{s_i} =
    1 + \frac{\delta_i}{\sqrt{|K|}} \cdot \sqrt{\frac{1}{p_i} - 1}.
\]

Under uniform key access the node with the biggest slowdown factor will
bottleneck the entire system, so we can define the cluster-wide slowdown factor as:
\begin{equation}
    \label{eq:slowdown}
    f = \max_i f_i = 1 + \frac{1}{\sqrt{|K|}} \cdot \max_i \delta_i \sqrt{\frac{1}{p_i} - 1}.
\end{equation}

When strict guarantees are required, the client can choose to overbook the
table's throughput to accommodate potential slowdowns.
